% Generated draft: requires mathematical and source review.
\documentclass[11pt]{article}
\usepackage[margin=25mm]{geometry}
\usepackage{fontspec}
\setmainfont{Latin Modern Roman}
\usepackage{amsmath,amssymb,mathtools,unicode-math}
\setmathfont{Latin Modern Math}
\usepackage{enumitem,needspace,xurl,hyperref}
\setlength{\parindent}{0pt}
\setlength{\parskip}{5pt}
\setlength{\emergencystretch}{3em}
\newcommand{\N}{\mathbb{N}}
\newcommand{\Z}{\mathbb{Z}}
\newcommand{\Q}{\mathbb{Q}}
\newcommand{\R}{\mathbb{R}}
\newcommand{\C}{\mathbb{C}}
\newcommand{\F}{\mathbb{F}}
\newcommand{\A}{\mathbb{A}}
\newcommand{\PP}{\mathbb{P}}
\newcommand{\E}{\mathbb{E}}
\newcommand{\eps}{\varepsilon}
\newcommand{\dd}{\,\mathrm{d}}
\begin{document}
\section*{UnsolvedMath ID: 30005343}
\textbf{Source problem code:} catalogue entry
\subsection{Definitions and mathematical statement}
Let \(n\in\Z_{\geq 1}\) and write \([n]=\{1,\ldots,n\}\). A Latin square of order \(n\) is a map \(L:[n]\times[n]\to[n]\) such that each symbol occurs exactly once in each row and exactly once in each column. A partial transversal is a set \(T\subseteq[n]\times[n]\) for which the three maps
\[
(i,j)\longmapsto i,\qquad (i,j)\longmapsto j,\qquad (i,j)\longmapsto L(i,j)
\]
are all injective on \(T\). Its size is \(|T|\). A full transversal is a partial transversal of size \(n\).

The supplied Ryser--Brualdi--Stein conjecture asks whether, for every positive integer \(n\) and every Latin square \(L\) of order \(n\), both of the following hold:
\begin{enumerate}
\item There exists a partial transversal \(T\) with \(|T|=n-1\).
\item If \(n\) is odd, there exists a full transversal \(U\) with \(|U|=n\).
\end{enumerate}
For \(n=1\), the first assertion permits the empty partial transversal. No full-transversal assertion is imposed for even \(n\).

The source's phrase ``transversal of size \(n-1\)'' is interpreted here as a partial transversal. The supplied original report excerpt is truncated and contains corrupted mathematical notation; it does not display the conjecture itself. This formulation therefore follows the supplied clean statement, whose report-verification label is a supplied curation claim and has not been independently checked.
\subsection{Short English statement}
Does every Latin square of positive order \(n\) contain \(n-1\) entries in distinct rows and columns with distinct symbols, and, whenever \(n\) is odd, \(n\) such entries? The registry labels the problem partially solved; that status is supplied dated curation, not an independently verified assessment.
\subsection{Sources}
\begin{itemize}
\item Local UnsolvedMath registry, permanent ID \(30005343\), source problem code \texttt{catalogue entry}, titled ``Ryser--Brualdi--Stein Transversal Conjecture,'' in the Oberwolfach Reports Open Problems collection. The record gives proposed year 2023 and names no individual proposer. Its \texttt{corrected\_verified} label and claim of checking against the official report are supplied metadata, not independent verification here.
\item Combinatorics (2023), Oberwolfach Reports, pp.\ 5--89, DOI 10.4171/owr/2023/1. Also listed in the supplied literature sources as ``Oberwolfach Reports 2023/1'': \url{https://doi.org/10.4171/owr/2023/1}.
\item Montgomery, ``Transversals in Latin Squares,'' supplied literature reference: \url{https://arxiv.org/abs/2310.19779}.
\item Montgomery, ``Latin Square Transversal Survey,'' supplied literature reference: \url{https://arxiv.org/abs/2406.19873}.
\item Supplied literature assessment dated exactly \texttt{2026-08-22}: the record reports that Montgomery proves the \(n-1\) partial-transversal conclusion for all sufficiently large orders, while the full Ryser--Brualdi--Stein assertion, including the odd-order full-transversal case, remains unresolved. This is supplied dated curation derived from the supplied Oberwolfach Reports curation, not an independently verified literature assessment. The record expressly requests independent verification of time-sensitive status and cited claims.
\end{itemize}
\end{document}
